\section{HTML5表单与多媒体元素}\label{html5ux8868ux5355ux4e0eux591aux5a92ux4f53ux5143ux7d20}

HTML5引入了多种新型表单输入类型、属性和多媒体元素，极大地提高了数据采集的效率和多媒体内容的呈现能力。在智慧水利平台中，这些功能对于数据录入和音视频监控尤为重要。

\subsection{1.
HTML5表单新特性}\label{html5ux8868ux5355ux65b0ux7279ux6027}

HTML5大幅增强了表单功能，通过新的输入类型、验证属性和UI优化，让用户可以更高效地输入数据。

\subsubsection{1.1
新的输入类型}\label{ux65b0ux7684ux8f93ux5165ux7c7bux578b}

\begin{lstlisting}[language=HTML]
<form id="water-data-form">
    <!-- 日期和时间类型 -->
    <div class="form-group">
        <label for="reading-date">观测日期：</label>
        <input type="date" id="reading-date" name="reading-date" required>
    </div>
    
    <div class="form-group">
        <label for="reading-time">观测时间：</label>
        <input type="time" id="reading-time" name="reading-time" required>
    </div>
    
    <div class="form-group">
        <label for="reading-datetime">完整时间戳：</label>
        <input type="datetime-local" id="reading-datetime" name="reading-datetime">
    </div>
    
    <!-- 数值类型 -->
    <div class="form-group">
        <label for="water-level">水位(米)：</label>
        <input type="number" id="water-level" name="water-level" 
               step="0.01" min="0" max="200" required>
    </div>
    
    <div class="form-group">
        <label for="flow-rate">流量设定：</label>
        <input type="range" id="flow-rate" name="flow-rate" 
               min="0" max="1000" step="10" value="500">
        <output for="flow-rate">500</output> 立方米/秒
    </div>
    
    <!-- 特殊输入类型 -->
    <div class="form-group">
        <label for="observer-email">观测员邮箱：</label>
        <input type="email" id="observer-email" name="observer-email" 
               placeholder="name@example.com">
    </div>
    
    <div class="form-group">
        <label for="station-url">站点网址：</label>
        <input type="url" id="station-url" name="station-url" 
               placeholder="https://example.com/station">
    </div>
    
    <div class="form-group">
        <label for="station-search">站点搜索：</label>
        <input type="search" id="station-search" name="station-search" 
               placeholder="输入站点编号或名称">
    </div>
    
    <div class="form-group">
        <label for="station-phone">联系电话：</label>
        <input type="tel" id="station-phone" name="station-phone" 
               pattern="[0-9]{3}-[0-9]{8}" placeholder="010-12345678">
    </div>
    
    <div class="form-group">
        <label for="water-color">水色选择：</label>
        <input type="color" id="water-color" name="water-color" 
               value="#a6d1f5">
    </div>
</form>
\end{lstlisting}

输入类型说明：

\begin{itemize}
\tightlist
\item
  \passthrough{\lstinline!date!}、\passthrough{\lstinline!time!}、\passthrough{\lstinline!datetime-local!}：日期和时间选择器
\item
  \passthrough{\lstinline!number!}：数字输入框，支持最小/最大值和步进值
\item
  \passthrough{\lstinline!range!}：滑块控件，适合快速选择范围值
\item
  \passthrough{\lstinline!email!}、\passthrough{\lstinline!url!}、\passthrough{\lstinline!tel!}：专用于特定格式的输入，带有相应验证
\item
  \passthrough{\lstinline!search!}：搜索框，通常带有清除按钮
\item
  \passthrough{\lstinline!color!}：颜色选择器
\end{itemize}

\subsubsection{1.2
新的表单验证属性}\label{ux65b0ux7684ux8868ux5355ux9a8cux8bc1ux5c5eux6027}

HTML5引入了内建的表单验证机制，无需JavaScript即可完成基本验证：

\begin{lstlisting}[language=HTML]
<form id="water-quality-form" novalidate>
    <div class="form-group">
        <label for="station-id">监测站点ID：</label>
        <input type="text" id="station-id" name="station-id" required 
               pattern="[A-Z]{2}[0-9]{4}" 
               title="格式：两位大写字母和四位数字，如AB1234">
    </div>
    
    <div class="form-group">
        <label for="ph-value">pH值：</label>
        <input type="number" id="ph-value" name="ph-value" 
               step="0.1" min="0" max="14" required>
    </div>
    
    <div class="form-group">
        <label for="dissolved-oxygen">溶解氧(mg/L)：</label>
        <input type="number" id="dissolved-oxygen" name="dissolved-oxygen" 
               step="0.01" min="0" max="20" required>
    </div>
    
    <div class="form-group">
        <label for="water-temperature">水温(°C)：</label>
        <input type="number" id="water-temperature" name="water-temperature" 
               step="0.1" min="-10" max="50">
    </div>
    
    <div class="form-group">
        <label for="notes">备注：</label>
        <textarea id="notes" name="notes" rows="3" maxlength="500" 
                  placeholder="输入观测备注信息..."></textarea>
    </div>
    
    <div class="form-actions">
        <button type="submit">提交数据</button>
        <button type="reset">重置</button>
    </div>
</form>
\end{lstlisting}

验证属性说明：

\begin{itemize}
\tightlist
\item
  \passthrough{\lstinline!required!}：指定必填字段
\item
  \passthrough{\lstinline!pattern!}：使用正则表达式定义输入格式
\item
  \passthrough{\lstinline!min!}/\passthrough{\lstinline!max!}：设定数值或日期的范围
\item
  \passthrough{\lstinline!step!}：定义数值的步进值
\item
  \passthrough{\lstinline!maxlength!}/\passthrough{\lstinline!minlength!}：限制文本长度
\item
  \passthrough{\lstinline!title!}：提供验证失败时的提示信息
\item
  \passthrough{\lstinline!novalidate!}：禁用表单的自动验证（通常配合JavaScript手动验证）
\end{itemize}

\subsubsection{1.3
表单元素和属性}\label{ux8868ux5355ux5143ux7d20ux548cux5c5eux6027}

HTML5引入了更多表单元素和属性，增强用户体验：

\begin{lstlisting}[language=HTML]
<form id="reservoir-data-form">
    <!-- 数据列表 -->
    <div class="form-group">
        <label for="reservoir-name">水库名称：</label>
        <input type="text" id="reservoir-name" name="reservoir-name" 
               list="reservoir-list" placeholder="选择或输入水库名称">
        <datalist id="reservoir-list">
            <option value="龙泉水库">
            <option value="明月湖">
            <option value="青山水库">
            <option value="长安水库">
            <option value="白云水库">
        </datalist>
    </div>
    
    <!-- 输出元素 -->
    <div class="form-group">
        <label for="storage-percent">蓄水比例：</label>
        <input type="range" id="storage-percent" name="storage-percent" 
               min="0" max="100" value="75" 
               oninput="storage_output.value = storage_percent.value + '%'">
        <output name="storage_output" id="storage_output" for="storage-percent">75%</output>
    </div>
    
    <!-- 进度条 -->
    <div class="form-group">
        <label>数据上传进度：</label>
        <progress id="upload-progress" max="100" value="0">0%</progress>
    </div>
    
    <!-- 仪表盘 -->
    <div class="form-group">
        <label for="water-quality">水质评级：</label>
        <meter id="water-quality" min="0" max="100" low="33" high="66" optimum="80" value="75">75/100</meter>
    </div>
    
    <!-- 自动完成 -->
    <div class="form-group">
        <label for="observer-name">观测员：</label>
        <input type="text" id="observer-name" name="observer-name" 
               autocomplete="name" placeholder="输入姓名">
    </div>
    
    <!-- 自动聚焦 -->
    <div class="form-group">
        <label for="reading-value">读数值：</label>
        <input type="number" id="reading-value" name="reading-value" autofocus>
    </div>
    
    <!-- 占位符文本 -->
    <div class="form-group">
        <label for="data-comments">备注信息：</label>
        <textarea id="data-comments" name="data-comments" 
                  placeholder="请输入观测过程中的特殊情况..."></textarea>
    </div>
</form>
\end{lstlisting}

新元素和属性说明：

\begin{itemize}
\tightlist
\item
  \passthrough{\lstinline!<datalist>!}：为输入框提供预定义选项
\item
  \passthrough{\lstinline!<output>!}：显示计算或操作的结果
\item
  \passthrough{\lstinline!<progress>!}：表示任务完成进度
\item
  \passthrough{\lstinline!<meter>!}：表示已知范围内的标量值
\item
  \passthrough{\lstinline!autocomplete!}：启用自动完成功能
\item
  \passthrough{\lstinline!autofocus!}：页面加载时自动获得焦点
\item
  \passthrough{\lstinline!placeholder!}：提供输入提示文本
\end{itemize}

\subsubsection{1.4
表单样式化和交互}\label{ux8868ux5355ux6837ux5f0fux5316ux548cux4ea4ux4e92}

HTML5表单可以通过CSS进行样式化，并结合JavaScript提供更丰富的交互：

\begin{lstlisting}[language=HTML]
<style>
    /* 基本样式 */
    .form-group {
        margin-bottom: 15px;
    }
    
    label {
        display: block;
        margin-bottom: 5px;
        font-weight: bold;
    }
    
    input, select, textarea {
        width: 100%;
        padding: 8px;
        border: 1px solid #ddd;
        border-radius: 4px;
    }
    
    /* 验证样式 */
    input:valid {
        border-color: #4CAF50;
    }
    
    input:invalid {
        border-color: #F44336;
    }
    
    input:focus {
        outline: none;
        box-shadow: 0 0 5px rgba(0,123,255,0.5);
    }
    
    /* 自定义范围滑块 */
    input[type="range"] {
        height: 8px;
        background: #e9ecef;
    }
    
    /* 自定义复选框 */
    input[type="checkbox"] {
        width: auto;
        margin-right: 8px;
    }
</style>

<script>
    // 表单验证示例
    document.getElementById('water-data-form').addEventListener('submit', function(event) {
        const waterLevel = document.getElementById('water-level').value;
        if (waterLevel > 180) {
            alert('警告：水位值异常，请确认！');
            event.preventDefault();
        }
    });
    
    // 范围控件与输出同步
    document.getElementById('flow-rate').addEventListener('input', function() {
        document.querySelector('output[for="flow-rate"]').textContent = this.value;
    });
</script>
\end{lstlisting}

\subsection{2.
HTML5多媒体元素}\label{html5ux591aux5a92ux4f53ux5143ux7d20}

HTML5引入的原生音视频支持，使得在网页中嵌入和控制多媒体内容变得简单，这对于水利监控系统至关重要。

\subsubsection{2.1 视频元素}\label{ux89c6ux9891ux5143ux7d20}

\begin{lstlisting}[language=HTML]
<div class="monitoring-video">
    <h3>大坝实时监控</h3>
    <video width="640" height="360" controls autoplay muted>
        <source src="/videos/dam-monitoring.mp4" type="video/mp4">
        <source src="/videos/dam-monitoring.webm" type="video/webm">
        <p>您的浏览器不支持HTML5视频播放。请升级或使用其他浏览器。</p>
    </video>
    
    <div class="video-controls">
        <button id="fullscreen-btn">全屏</button>
        <button id="snapshot-btn">截图</button>
        <select id="camera-selector">
            <option value="camera1">摄像头1：大坝上游</option>
            <option value="camera2">摄像头2：溢洪道</option>
            <option value="camera3">摄像头3：下游河道</option>
        </select>
    </div>
</div>
\end{lstlisting}

视频元素属性：

\begin{itemize}
\tightlist
\item
  \passthrough{\lstinline!controls!}：显示视频控制界面
\item
  \passthrough{\lstinline!autoplay!}：自动开始播放
\item
  \passthrough{\lstinline!muted!}：静音播放
\item
  \passthrough{\lstinline!loop!}：循环播放
\item
  \passthrough{\lstinline!poster!}：在视频加载前显示的图像
\item
  \passthrough{\lstinline!preload!}：预加载策略（\passthrough{\lstinline!auto!},
  \passthrough{\lstinline!metadata!}, \passthrough{\lstinline!none!}）
\end{itemize}

\subsubsection{2.2 音频元素}\label{ux97f3ux9891ux5143ux7d20}

\begin{lstlisting}[language=HTML]
<div class="audio-alert">
    <h3>警报声音系统</h3>
    <audio id="flood-warning" controls>
        <source src="/audio/flood-warning.mp3" type="audio/mpeg">
        <source src="/audio/flood-warning.ogg" type="audio/ogg">
        <p>您的浏览器不支持HTML5音频播放。</p>
    </audio>
    
    <button id="play-warning">播放警报</button>
    <button id="stop-warning">停止警报</button>
    
    <div class="volume-control">
        <label for="volume">音量：</label>
        <input type="range" id="volume" min="0" max="1" step="0.1" value="0.7">
    </div>
</div>

<script>
    const audio = document.getElementById('flood-warning');
    const volumeControl = document.getElementById('volume');
    
    document.getElementById('play-warning').addEventListener('click', function() {
        audio.play();
    });
    
    document.getElementById('stop-warning').addEventListener('click', function() {
        audio.pause();
        audio.currentTime = 0;
    });
    
    volumeControl.addEventListener('input', function() {
        audio.volume = this.value;
    });
</script>
\end{lstlisting}

\subsubsection{2.3 图形与绘图}\label{ux56feux5f62ux4e0eux7ed8ux56fe}

HTML5提供了强大的绘图工具，适用于水利数据可视化：

\paragraph{2.3.1 Canvas元素}\label{canvasux5143ux7d20}

\begin{lstlisting}[language=HTML]
<div class="water-level-chart">
    <h3>近24小时水位变化</h3>
    <canvas id="waterLevelCanvas" width="800" height="400"></canvas>
</div>

<script>
    const canvas = document.getElementById('waterLevelCanvas');
    const ctx = canvas.getContext('2d');
    
    // 模拟数据
    const times = ['00:00', '02:00', '04:00', '06:00', '08:00', '10:00', 
                  '12:00', '14:00', '16:00', '18:00', '20:00', '22:00'];
    const levels = [142.5, 142.8, 143.2, 144.1, 145.0, 145.3, 
                    145.2, 144.9, 144.5, 144.0, 143.5, 143.0];
    
    // 设置样式
    ctx.fillStyle = '#f5f5f5';
    ctx.fillRect(0, 0, canvas.width, canvas.height);
    
    // 绘制坐标轴
    ctx.beginPath();
    ctx.moveTo(50, 50);
    ctx.lineTo(50, 350);
    ctx.lineTo(750, 350);
    ctx.strokeStyle = '#333';
    ctx.lineWidth = 2;
    ctx.stroke();
    
    // 绘制数据点和线
    ctx.beginPath();
    for (let i = 0; i < times.length; i++) {
        const x = 50 + i * (700 / (times.length - 1));
        const y = 350 - ((levels[i] - 142) / 4 * 300);
        
        if (i === 0) {
            ctx.moveTo(x, y);
        } else {
            ctx.lineTo(x, y);
        }
        
        // 绘制数据点
        ctx.fillStyle = '#1890ff';
        ctx.beginPath();
        ctx.arc(x, y, 5, 0, Math.PI * 2);
        ctx.fill();
        
        // 绘制坐标标签
        ctx.fillStyle = '#666';
        ctx.font = '12px Arial';
        ctx.textAlign = 'center';
        ctx.fillText(times[i], x, 370);
        ctx.fillText(levels[i].toFixed(1) + 'm', x, y - 15);
    }
    
    // 设置线样式并描边
    ctx.strokeStyle = '#1890ff';
    ctx.lineWidth = 3;
    ctx.stroke();
    
    // 绘制警戒水位线
    ctx.beginPath();
    const warningLevel = 145;
    const warningY = 350 - ((warningLevel - 142) / 4 * 300);
    ctx.moveTo(50, warningY);
    ctx.lineTo(750, warningY);
    ctx.setLineDash([5, 5]);
    ctx.strokeStyle = '#ff4d4f';
    ctx.lineWidth = 2;
    ctx.stroke();
    
    // 添加警戒线标签
    ctx.fillStyle = '#ff4d4f';
    ctx.font = '14px Arial';
    ctx.textAlign = 'left';
    ctx.fillText('警戒水位: ' + warningLevel + 'm', 60, warningY - 10);
</script>
\end{lstlisting}

\paragraph{2.3.2 SVG向量图形}\label{svgux5411ux91cfux56feux5f62}

\begin{lstlisting}[language=HTML]
<div class="reservoir-diagram">
    <h3>水库示意图</h3>
    <svg width="800" height="400" viewBox="0 0 800 400">
        <!-- 背景 -->
        <rect x="0" y="0" width="800" height="400" fill="#f0f9ff" />
        
        <!-- 水库轮廓 -->
        <path d="M100,100 C200,80 400,50 700,100 L700,300 C400,350 200,320 100,300 Z" 
              fill="#a6d1f5" stroke="#0088d1" stroke-width="2" />
        
        <!-- 大坝 -->
        <rect x="50" y="150" width="50" height="150" fill="#8d6e63" />
        
        <!-- 水流动画 -->
        <path d="M50,220 C20,220 20,230 50,230" stroke="#5aadff" stroke-width="5" 
              stroke-linecap="round">
            <animate attributeName="d" 
                     values="M50,220 C20,220 20,230 50,230; 
                             M50,220 C40,225 30,225 50,230; 
                             M50,220 C20,220 20,230 50,230" 
                     dur="2s" repeatCount="indefinite" />
        </path>
        
        <!-- 监测点 -->
        <g class="monitoring-point" transform="translate(300, 150)">
            <circle cx="0" cy="0" r="8" fill="#ff7a45" />
            <circle cx="0" cy="0" r="12" fill="none" stroke="#ff7a45" stroke-width="2">
                <animate attributeName="r" values="12;20;12" dur="3s" repeatCount="indefinite" />
                <animate attributeName="opacity" values="1;0;1" dur="3s" repeatCount="indefinite" />
            </circle>
            <text x="15" y="5" fill="#333">监测点A</text>
        </g>
        
        <g class="monitoring-point" transform="translate(500, 200)">
            <circle cx="0" cy="0" r="8" fill="#597ef7" />
            <circle cx="0" cy="0" r="12" fill="none" stroke="#597ef7" stroke-width="2">
                <animate attributeName="r" values="12;20;12" dur="3s" repeatCount="indefinite" />
                <animate attributeName="opacity" values="1;0;1" dur="3s" repeatCount="indefinite" />
            </circle>
            <text x="15" y="5" fill="#333">监测点B</text>
        </g>
        
        <!-- 标注 -->
        <text x="400" y="30" font-size="18" text-anchor="middle" fill="#333">龙泉水库示意图</text>
        <text x="150" y="350" font-size="14" fill="#333">水位: 145.6米</text>
        <text x="600" y="350" font-size="14" fill="#333">蓄水量: 2.3亿立方米</text>
    </svg>
</div>
\end{lstlisting}

\subsection{3.
在智慧水利平台中的应用}\label{ux5728ux667aux6167ux6c34ux5229ux5e73ux53f0ux4e2dux7684ux5e94ux7528}

HTML5表单和多媒体元素在智慧水利平台中有广泛应用，以下是几个具体场景：

\subsubsection{3.1
水文数据采集表单}\label{ux6c34ux6587ux6570ux636eux91c7ux96c6ux8868ux5355}

\begin{lstlisting}[language=HTML]
<div class="data-collection-module">
    <h2>水文数据采集</h2>
    <form id="hydrology-data-form" class="smart-form">
        <fieldset>
            <legend>基本信息</legend>
            <div class="form-row">
                <div class="form-group">
                    <label for="station-code">测站编码</label>
                    <input type="text" id="station-code" name="station-code" required 
                           pattern="[A-Z]{2}[0-9]{5}" 
                           title="格式：两位大写字母+五位数字">
                </div>
                <div class="form-group">
                    <label for="collection-time">采集时间</label>
                    <input type="datetime-local" id="collection-time" name="collection-time" required>
                </div>
                <div class="form-group">
                    <label for="collection-type">采集类型</label>
                    <select id="collection-type" name="collection-type" required>
                        <option value="">--请选择--</option>
                        <option value="regular">常规观测</option>
                        <option value="emergency">应急监测</option>
                        <option value="scheduled">定时监测</option>
                    </select>
                </div>
            </div>
        </fieldset>
        
        <fieldset>
            <legend>水文参数</legend>
            <div class="form-row">
                <div class="form-group">
                    <label for="water-level">水位(m)</label>
                    <input type="number" id="water-level" name="water-level" 
                           step="0.01" min="0" max="200">
                </div>
                <div class="form-group">
                    <label for="flow-rate">流量(m³/s)</label>
                    <input type="number" id="flow-rate" name="flow-rate" 
                           step="0.1" min="0">
                </div>
                <div class="form-group">
                    <label for="flow-velocity">流速(m/s)</label>
                    <input type="number" id="flow-velocity" name="flow-velocity" 
                           step="0.01" min="0">
                </div>
            </div>
        </fieldset>
        
        <fieldset>
            <legend>水质参数</legend>
            <div class="form-row">
                <div class="form-group">
                    <label for="temperature">水温(°C)</label>
                    <input type="number" id="temperature" name="temperature" 
                           step="0.1" min="-10" max="50">
                </div>
                <div class="form-group">
                    <label for="ph">pH值</label>
                    <input type="number" id="ph" name="ph" 
                           step="0.1" min="0" max="14">
                    <meter value="7.2" min="0" max="14" low="6" high="9" optimum="7">7.2</meter>
                </div>
                <div class="form-group">
                    <label for="turbidity">浊度(NTU)</label>
                    <input type="number" id="turbidity" name="turbidity" 
                           step="0.1" min="0">
                </div>
            </div>
        </fieldset>
        
        <fieldset>
            <legend>附加信息</legend>
            <div class="form-group">
                <label for="weather">天气状况</label>
                <input type="text" id="weather" name="weather" 
                       list="weather-conditions">
                <datalist id="weather-conditions">
                    <option value="晴朗">
                    <option value="多云">
                    <option value="小雨">
                    <option value="中雨">
                    <option value="大雨">
                    <option value="暴雨">
                </datalist>
            </div>
            <div class="form-group">
                <label for="field-photo">现场照片</label>
                <input type="file" id="field-photo" name="field-photo" 
                       accept="image/*" capture="environment">
                <div class="preview-area" id="photo-preview"></div>
            </div>
            <div class="form-group">
                <label for="notes">观测备注</label>
                <textarea id="notes" name="notes" rows="3"></textarea>
            </div>
        </fieldset>
        
        <div class="form-actions">
            <button type="submit" class="btn-primary">提交数据</button>
            <button type="reset" class="btn-secondary">重置表单</button>
            <button type="button" class="btn-outline" id="save-draft">保存草稿</button>
        </div>
    </form>
</div>
\end{lstlisting}

\subsubsection{3.2
视频监控系统}\label{ux89c6ux9891ux76d1ux63a7ux7cfbux7edf}

\begin{lstlisting}[language=HTML]
<div class="monitoring-system">
    <h2>水库视频监控系统</h2>
    
    <div class="video-grid">
        <div class="video-card primary">
            <h3>主坝监控</h3>
            <video id="main-dam-video" width="640" height="360" autoplay muted>
                <source src="/streams/main-dam-live.m3u8" type="application/x-mpegURL">
                <p>您的浏览器不支持HTML5视频播放。</p>
            </video>
            <div class="video-controls">
                <button id="toggle-play" class="icon-button" title="播放/暂停">
                    <i class="icon-pause"></i>
                </button>
                <button id="toggle-mute" class="icon-button" title="静音/取消静音">
                    <i class="icon-mute"></i>
                </button>
                <button id="take-snapshot" class="icon-button" title="截图">
                    <i class="icon-camera"></i>
                </button>
                <button id="toggle-fullscreen" class="icon-button" title="全屏">
                    <i class="icon-fullscreen"></i>
                </button>
                <div class="ptz-controls">
                    <button id="ptz-up" class="icon-button" title="向上">⬆</button>
                    <button id="ptz-down" class="icon-button" title="向下">⬇</button>
                    <button id="ptz-left" class="icon-button" title="向左">⬅</button>
                    <button id="ptz-right" class="icon-button" title="向右">➡</button>
                    <button id="ptz-zoomin" class="icon-button" title="放大">+</button>
                    <button id="ptz-zoomout" class="icon-button" title="缩小">-</button>
                </div>
            </div>
            <div class="video-info">
                <div class="camera-info">
                    <span class="label">摄像头ID：</span>
                    <span class="value">CAM001</span>
                </div>
                <div class="location-info">
                    <span class="label">位置：</span>
                    <span class="value">主坝中段</span>
                </div>
                <div class="status-info">
                    <span class="label">状态：</span>
                    <span class="value online">在线</span>
                </div>
                <div class="time-info">
                    <span class="label">时间：</span>
                    <span class="value" id="live-time">2023-06-15 15:30:22</span>
                </div>
            </div>
        </div>
        
        <div class="thumbnail-grid">
            <div class="video-thumbnail" data-camera-id="CAM002">
                <video width="160" height="90" muted loop>
                    <source src="/streams/spillway-preview.mp4" type="video/mp4">
                </video>
                <div class="thumbnail-overlay">
                    <span class="camera-name">溢洪道</span>
                </div>
            </div>
            
            <div class="video-thumbnail" data-camera-id="CAM003">
                <video width="160" height="90" muted loop>
                    <source src="/streams/power-station-preview.mp4" type="video/mp4">
                </video>
                <div class="thumbnail-overlay">
                    <span class="camera-name">发电站</span>
                </div>
            </div>
            
            <div class="video-thumbnail" data-camera-id="CAM004">
                <video width="160" height="90" muted loop>
                    <source src="/streams/upstream-preview.mp4" type="video/mp4">
                </video>
                <div class="thumbnail-overlay">
                    <span class="camera-name">上游河道</span>
                </div>
            </div>
            
            <div class="video-thumbnail" data-camera-id="CAM005">
                <video width="160" height="90" muted loop>
                    <source src="/streams/downstream-preview.mp4" type="video/mp4">
                </video>
                <div class="thumbnail-overlay">
                    <span class="camera-name">下游河道</span>
                </div>
            </div>
        </div>
    </div>
    
    <div class="snapshots-container">
        <h3>历史截图</h3>
        <div class="snapshots-grid" id="snapshots-grid">
            <!-- 动态加载的截图将显示在这里 -->
        </div>
    </div>
</div>

<script>
    // 视频控制功能示例
    document.getElementById('toggle-play').addEventListener('click', function() {
        const video = document.getElementById('main-dam-video');
        if (video.paused) {
            video.play();
            this.querySelector('i').className = 'icon-pause';
        } else {
            video.pause();
            this.querySelector('i').className = 'icon-play';
        }
    });
    
    document.getElementById('take-snapshot').addEventListener('click', function() {
        const video = document.getElementById('main-dam-video');
        const canvas = document.createElement('canvas');
        canvas.width = video.videoWidth;
        canvas.height = video.videoHeight;
        
        const ctx = canvas.getContext('2d');
        ctx.drawImage(video, 0, 0, canvas.width, canvas.height);
        
        // 创建截图缩略图
        const snapshot = document.createElement('div');
        snapshot.className = 'snapshot-item';
        
        const img = document.createElement('img');
        img.src = canvas.toDataURL('image/png');
        
        const timeLabel = document.createElement('div');
        timeLabel.className = 'time-label';
        timeLabel.textContent = new Date().toLocaleString();
        
        snapshot.appendChild(img);
        snapshot.appendChild(timeLabel);
        
        // 添加到截图网格
        document.getElementById('snapshots-grid').prepend(snapshot);
    });
    
    // 更新实时时间
    setInterval(function() {
        document.getElementById('live-time').textContent = 
            new Date().toLocaleString('zh-CN');
    }, 1000);
    
    // 缩略图点击事件
    document.querySelectorAll('.video-thumbnail').forEach(function(thumbnail) {
        thumbnail.addEventListener('click', function() {
            // 在实际应用中，这里会切换主视频源
            alert('切换到摄像头: ' + this.dataset.cameraId);
        });
    });
</script>
\end{lstlisting}

\subsection{4.
最佳实践与注意事项}\label{ux6700ux4f73ux5b9eux8df5ux4e0eux6ce8ux610fux4e8bux9879}

在智慧水利平台开发中使用HTML5表单和多媒体元素时，应注意以下几点：

\subsubsection{4.1
表单设计原则}\label{ux8868ux5355ux8bbeux8ba1ux539fux5219}

\begin{enumerate}
\def\labelenumi{\arabic{enumi}.}
\tightlist
\item
  \textbf{用户友好性}：设计直观的表单布局，分组相关字段，使用合适的输入类型
\item
  \textbf{渐进增强}：确保基本功能在所有浏览器中可用，高级特性作为增强
\item
  \textbf{即时反馈}：提供清晰的验证反馈，帮助用户正确输入
\item
  \textbf{移动优化}：确保表单在移动设备上易于使用，考虑触摸交互
\item
  \textbf{数据安全}：敏感数据使用HTTPS传输，避免缓存敏感信息
\end{enumerate}

\subsubsection{4.2
多媒体性能优化}\label{ux591aux5a92ux4f53ux6027ux80fdux4f18ux5316}

\begin{enumerate}
\def\labelenumi{\arabic{enumi}.}
\tightlist
\item
  \textbf{响应式设计}：使用相对单位设置尺寸，确保在不同设备上正常显示
\item
  \textbf{懒加载}：对非关键视频使用懒加载，减少初始页面加载时间
\item
  \textbf{格式选择}：提供多种格式选项，如MP4/WebM，确保更广泛的兼容性
\item
  \textbf{带宽考虑}：提供不同分辨率的视频，根据网络条件自适应
\item
  \textbf{回退方案}：为不支持HTML5多媒体元素的用户提供替代方案
\end{enumerate}

\subsubsection{4.3
辅助功能与可访问性}\label{ux8f85ux52a9ux529fux80fdux4e0eux53efux8bbfux95eeux6027}

\begin{enumerate}
\def\labelenumi{\arabic{enumi}.}
\tightlist
\item
  \textbf{标签关联}：确保每个表单控件都有关联的label
\item
  \textbf{键盘导航}：确保所有交互元素可通过键盘访问
\item
  \textbf{多媒体替代文本}：为视频提供字幕和文字说明
\item
  \textbf{ARIA属性}：使用适当的ARIA角色和属性增强可访问性
\item
  \textbf{颜色对比度}：确保文本和控件有足够的对比度
\end{enumerate}

\subsection{5. 总结}\label{ux603bux7ed3}

HTML5表单与多媒体元素为智慧水利平台提供了丰富的用户交互和内容呈现能力。通过表单新特性，可以更高效地收集和验证水利监测数据；通过多媒体元素，可以直观地展示水利工程的实时状况。这些技术的合理应用将大幅提升平台的用户体验，增强数据可视化效果，为水利工程的智能化监测与管理提供有力的技术支持。
